% SEGMENT OLP-0101-S01
% Part: first-order-logic
% Chapter: tableaux
% Section: quantifier-rules

\documentclass[../../../include/open-logic-section]{subfiles}

\begin{document}

\olfileid{fol}{tab}{qrl}

\olsection{அளவையடை விதிகள்}

\subsection{$\lforall$ க்கான விதிகள்}

\begin{defish}
\AxiomC{\sFmla{\True}{\lforall[x][!A(x)]}}
% SOURCE CORRECTION TA-TAB-002: use the edition's logical quantifier macro.
\RightLabel{\TRule{\True}{\lforall}}
\UnaryInfC{\sFmla{\True}{!A(t)}}
\DisplayProof
\hfill
\AxiomC{\sFmla{\False}{\lforall[x][!A(x)]}}
\RightLabel{\TRule{\False}{\lforall}}
\UnaryInfC{\sFmla{\False}{!A(a)}}
\DisplayProof
\end{defish}

\TRule{\True}{\lforall} இல், $t$ ஒரு மூடிய சொல் (அதாவது மாறிகள் இல்லாத
சொல்). \TRule{\False}{\lforall} இல், $a$ என்பது
\TRule{\False}{\lforall} விதிக்கு மேலுள்ள கிளையில் எங்கும் இடம்பெறாத
!!a{constant}. $a$ ஐ \TRule{\False}{\lforall} அனுமானத்தின்
\emph{தனிமாறி} என்கிறோம்.\footnote{மேலுள்ள விதியில் $a$ !!a{constant}
ஆக இருந்தாலும் ``தனிமாறி'' என்ற சொல்லையே பயன்படுத்துகிறோம். இதற்கு வரலாற்றுக்
காரணங்கள் உள்ளன.}

% SEGMENT OLP-0101-S02
\subsection{$\lexists$ க்கான விதிகள்}

\begin{defish}
\AxiomC{\sFmla{\True}{\lexists[x][!A(x)]}}
\RightLabel{\TRule{\True}{\lexists}}
\UnaryInfC{\sFmla{\True}{!A(a)}}
\DisplayProof
\hfill
\AxiomC{\sFmla{\False}{\lexists[x][!A(x)]}}
\RightLabel{\TRule{\False}{\lexists}}
\UnaryInfC{\sFmla{\False}{!A(t)}}
\DisplayProof
\end{defish}

மீண்டும், $t$ ஒரு மூடிய சொல்; $a$ என்பது \TRule{\True}{\lexists}
விதிக்கு மேலுள்ள கிளையில் இடம்பெறாத !!a{constant}. $a$ ஐ
\TRule{\True}{\lexists} அனுமானத்தின் \emph{தனிமாறி} என்கிறோம்.

% SEGMENT OLP-0101-S03
\TRule{\False}{\lforall} அல்லது \TRule{\True}{\lexists} அனுமானத்திற்கு
மேலுள்ள கிளையில் ஒரு தனிமாறி இடம்பெறக்கூடாது என்ற நிபந்தனை
\emph{தனிமாறி நிபந்தனை} எனப்படுகிறது.

% SEGMENT OLP-0101-S04
\begin{explain}
$!A$ என்பது !!{variable}~$x$ ஐக் கட்டற்றதாகக் கொண்ட !!a{formula} என்றால், அதை $!A(x)$
என்று எழுதிக் குறிப்போம் என்ற மரபை நினைவுகூர்க. அதே சூழலில் $!A(t)$ என்பது
$\Subst{!A}{t}{x}$ என்பதன் சுருக்கம். ஆகவே \TRule{\False}{\lexists}
விதியைப் பின்வருமாறும் எழுதலாம்:
\begin{prooftree}
  \AxiomC{\sFmla{\False}{\lexists[x][!A]}}
  \RightLabel{\TRule{\False}{\lexists}}
  \UnaryInfC{\sFmla{\False}{\Subst{!A}{t}{x}}}
\end{prooftree}
$t$ ஏற்கெனவே~$!A$ இல் இடம்பெறலாம் என்பதை கவனிக்கவும்; எடுத்துக்காட்டாக,
$!A$ ஆனது~$\Atom{\Obj P}{t,x}$ ஆக இருக்கலாம். எனவே
$\sFmla{\False}{\lexists[x][\Atom{\Obj P}{t,x}]}$ இலிருந்து
$\sFmla{\False}{\Atom{\Obj P}{t,t}}$ ஐ அனுமானிப்பது
\TRule{\False}{\lexists} இன் சரியான பயன்பாடாகும். ஆனால்
\TRule{\False}{\lforall}, \TRule{\True}{\lexists} ஆகியவற்றின்
தனிமாறி நிபந்தனைகள் !!{constant}~$a$ ஆனது~$!A$ இல் இடம்பெறக்கூடாது எனக்
கோருகின்றன. ஆகவே $\TRule{\False}{\lforall}$ ஐப் பயன்படுத்தி
$\sFmla{\False}{\lforall[x][\Atom{\Obj P}{a,x}]}$ இலிருந்து
$\sFmla{\False}{\Atom{\Obj P}{a,a}}$ ஐச் சரியாக அனுமானிக்க முடியாது.
\end{explain}

% SEGMENT OLP-0101-S05
\begin{explain}
\TRule{\True}{\lforall}, \TRule{\False}{\lexists} ஆகியவற்றில்
சொல்~$t$ மீது எந்தக் கட்டுப்பாடும் இல்லை. மறுபுறம்,
\TRule{\True}{\lexists}, \TRule{\False}{\lforall} விதிகளில்,
அந்தந்த அனுமானத்திற்கு மேலுள்ள கிளைகளில் !!{constant}~$a$ எங்கும்
இடம்பெறக்கூடாது என்று தனிமாறி நிபந்தனை கோருகிறது. முறைமை சரித்தன்மையுடையதாக
இருப்பதை உறுதிசெய்ய இது தேவை. இந்நிபந்தனை இல்லையெனில்,
$\lexists[x][\formula{A}(x)] \lif \lforall[x][\formula{A}(x)]$ க்குப்
பின்வரும் மூடிய !!a{tableau} கிடைக்கும்:
\begin{center}
\begin{tableau}{}
  [\sFmla{\False}{\lexists[x][\formula{A}(x)] \lif \lforall[x][\formula{A}(x)]}, just=\TAss
    [\sFmla{\True}{\lexists[x][\formula{A}(x)]},
      just={\TRule{\False}{\lif}[1]}
      [\sFmla{\False}{\lforall[x][\formula{A}(x)]},
        just={\TRule{\False}{\lif}[1]}
        [\sFmla{\True}{\formula{A}(a)},
          just={\TRule{\True}{\lexists}[2]}
          [\sFmla{\False}{\formula{A}(a)},
            just={\TRule{\False}{\lforall}[3]}, close]
        ]
      ]
    ]
  ]
\end{tableau}
\end{center}
ஆனால் $\lexists[x][\formula{A}(x)] \lif
\lforall[x][\formula{A}(x)]$ செல்லுபடியாகாது.
\end{explain}

\end{document}
